\section{工程师角色的重新定义}

\subsection{从写代码到设计环境}

传统软件工程中，工程师的核心产出是代码。在 Agent-First 模式下，这一角色发生了根本性转变。

\begin{importantbox}{工程师的新职责}
在 Agent-First 开发中，工程师的主要工作变为：
\begin{itemize}[nosep]
  \item \textbf{设计环境（Environment Design）}：构建 Agent 能有效工作的基础设施
  \item \textbf{表达意图（Intent Specification）}：将需求转化为 Agent 可理解的 Prompt
  \item \textbf{构建反馈循环（Feedback Loops）}：让 Agent 能自我验证和迭代
  \item \textbf{识别能力缺口}：当 Agent 失败时，分析缺少什么工具或抽象
\end{itemize}
\end{importantbox}

文章指出，早期进展比预期慢，\textbf{原因不是 Codex 能力不足，而是环境欠缺规范}（underspecified）。Agent 缺乏完成高层目标所需的工具、抽象和内部结构。因此，工程团队的首要任务变成了``赋能 Agent 做有用的工作''。

\subsection{深度优先的构建策略}

实践中，团队采用\textbf{深度优先}的构建策略：

\begin{enumerate}[nosep]
  \item 将大目标分解为更小的构建块（设计、编码、评审、测试等）
  \item 提示 Agent 构建这些基础模块
  \item 用这些模块解锁更复杂的任务
\end{enumerate}

当某个步骤失败时，修复方式几乎从来不是``再试一次''。因为唯一的推进方式是让 Codex 完成工作，人类工程师总是退一步问自己：\textbf{``缺少什么能力？如何让它对 Agent 既可见（legible）又可执行（enforceable）？''}

\begin{warningbox}{常见误区：把 Agent 当作更快的人类开发者}
很多团队在尝试 Agent 开发时，仍然用人类开发者的思维模式来使用 Agent。例如直接给一个模糊的大需求，期望 Agent 能像资深工程师一样自主分解和执行。文章的经验表明，这种方式会失败。Agent 需要的是\textbf{结构化的环境和明确的约束}，而不是更好的指令。``Try harder'' 几乎从来不是正确的解法。
\end{warningbox}

\subsection{人类与 Agent 的交互模式}

工程师几乎完全通过 Prompt 与系统交互。一个典型的工作流程如下：

\begin{enumerate}[nosep]
  \item 工程师描述一个任务
  \item 运行 Agent，让它开启一个 Pull Request
  \item 指示 Codex 在本地自我审查变更
  \item 请求其他 Agent 进行额外评审（本地和云端）
  \item 响应人类或 Agent 的反馈
  \item 在循环中迭代，直到所有 Agent 评审者满意
  \item 合并变更
\end{enumerate}

值得注意的是，人类\textbf{可以但不必须}审查 Pull Request。随着时间推移，团队将几乎所有审查工作推向了 Agent-to-Agent 的模式。

\begin{knowledgebox}{Ralph Wiggum Loop}
文章中提到了一个有趣的术语``Ralph Wiggum Loop''，指的是 Agent 自我审查、收到反馈、修改、再审查的循环迭代过程。这个名称来自《辛普森一家》中的角色 Ralph Wiggum，暗含一种自我对话的幽默感。这种循环是 Agent-First 开发中质量保障的核心机制。
\end{knowledgebox}

\subsection{本章小结}

在 Agent-First 范式下，工程师的角色从``代码编写者''变为``系统架构师 + Agent 教练''。核心技能从编程转变为环境设计、意图表达和反馈循环构建。关键原则是：当 Agent 失败时，不要``试更多次''，而是退一步问``缺少什么基础设施''。

\newpage

